\begin{helper}
Fix a terminal coordinate set \(U\), fixed blue and red present supports
\(B,R\subseteq U\) with \(|B|=u_B\), \(|R|=u_R\), and \(B\cap R=\emptyset\).
Let \(\tau=(P_\tau,A_\tau,Z_\tau)\) be a realized transcript for one fixed
\((B,J,\beta)\), and assume the RedBranch geometry
\[
P_\tau,A_\tau\subseteq U,\qquad P_\tau\cap A_\tau=\emptyset,\qquad
B\cup R\subseteq P_\tau,\qquad Z_\tau\subseteq A_\tau .
\]
Let \(\mathcal R_{\beta,\tau}\) be the released residual block for this
fixed realized pair.  Assume \(0\le p\le 1/2\), and assign complete residual
terminal assignments \(A\subseteq U\) the product mass
\[
\mu(A)=p^{|A|}(1-p)^{|U|-|A|}.
\]
Assume the RISI scan-stability conclusion for this fixed pair: every
residual assignment \(A\subseteq U\) satisfying
\[
P_\tau\setminus(B\cup R)\subseteq A,\qquad
A\cap(A_\tau\setminus Z_\tau)=\emptyset
\]
reconstructs to the same branch label \(\beta\) and the same transcript
\(\tau\).  Assume also that the released block is complete for this scan
output: every residual assignment with output \((\beta,\tau)\) belongs to
\(\mathcal R_{\beta,\tau}\).  If
\[
\rho_{\beta,\tau}
  =\sum_{A\in\mathcal R_{\beta,\tau}}\mu(A),
\]
then
\[
p^{|P_\tau|-u_R-u_B}(1-p)^{|A_\tau|-|Z_\tau|}
  \le \rho_{\beta,\tau}.
\]
\end{helper}
\begin{proof}
Let
\[
P^\circ=P_\tau\setminus(B\cup R),\qquad
A^\circ=A_\tau\setminus Z_\tau .
\]
The displayed RedBranch geometry gives \(P^\circ,A^\circ\subseteq U\).
Since \(P_\tau\cap A_\tau=\emptyset\), the smaller sets
\(P^\circ\) and \(A^\circ\) are disjoint.  Because \(B\cup R\subseteq
P_\tau\), \(B\cap R=\emptyset\), \(|B|=u_B\), and \(|R|=u_R\),
\[
|P^\circ|=|P_\tau|-u_R-u_B .
\]
Similarly, \(Z_\tau\subseteq A_\tau\) gives
\[
|A^\circ|=|A_\tau|-|Z_\tau|.
\]

Consider the residual cylinder
\[
\mathcal C_{\beta,\tau}
  =\{A\subseteq U:P^\circ\subseteq A,\ A\cap A^\circ=\emptyset\}.
\]
For each \(A\in\mathcal C_{\beta,\tau}\), the RISI scan-stability hypothesis
is precisely the fixed-\((B,J,\beta,\tau)\) consequence supplied by
\(\noderef{FixedSetHistoryCellRISIResidualScanStability}\): after the fixed
present support \(B\cup R\) is reinserted and the selected absences
\(Z_\tau\) are released, the reconstructed terminal truth table agrees with
the realized witness on every dependency set read by the reveal order, so the
same branch and transcript are returned.  The released-block completeness
hypothesis then puts this \(A\) in \(\mathcal R_{\beta,\tau}\).  Hence
\[
\mathcal C_{\beta,\tau}\subseteq\mathcal R_{\beta,\tau}.
\]
This is the point at which selected-sibling pruning is accounted for: a
residual assignment in the displayed cylinder is not discarded merely because
some selected coordinate was released; if its scan output is
\((\beta,\tau)\), completeness keeps it inside the corresponding released
block.

By the finite cylinder-mass identity
\(\noderef{FinsetPowersetCylinderMass}\), applied to the disjoint present
set \(P^\circ\) and absent set \(A^\circ\),
\[
\sum_{A\in\mathcal C_{\beta,\tau}}\mu(A)
  =p^{|P^\circ|}(1-p)^{|A^\circ|}
  =
p^{|P_\tau|-u_R-u_B}(1-p)^{|A_\tau|-|Z_\tau|}.
\]
The equality of the absent exponent is the same selected-absence factor split
recorded in \(\noderef{FixedSetHistoryCellSelectedAbsencePowSplit}\).  Since
\(0\le p\le1/2\), both \(p\) and \(1-p\) are nonnegative, so every
\(\mu(A)\) is nonnegative.  Therefore monotonicity of finite sums under the
inclusion \(\mathcal C_{\beta,\tau}\subseteq\mathcal R_{\beta,\tau}\) gives
\[
p^{|P_\tau|-u_R-u_B}(1-p)^{|A_\tau|-|Z_\tau|}
  =\sum_{A\in\mathcal C_{\beta,\tau}}\mu(A)
  \le
  \sum_{A\in\mathcal R_{\beta,\tau}}\mu(A)
  =\rho_{\beta,\tau}.
\]
\end{proof}
